\begin{figure}[t]
\centering
\begin{tikzpicture}[>=Latex,font=\small,x=1cm,y=1cm]
\fill[blue!8] (0,-.23) rectangle (3.3,.23);
\fill[blue!8] (3.0,-2.23) rectangle (6.3,-1.77);
\draw[->] (-.2,0)--(6.6,0) node[right] {$a$};
\draw[->] (-.2,-2)--(6.6,-2) node[right] {order in $z-i$};
\foreach \x/\lab in {0/0,3.3/A,6.3/{N-1}} {
 \draw (\x,-.08)--(\x,.08);
 \node[above=4pt] at (\x,0) {$\lab$};
}
\draw (0,-2.08)--(0,-1.92);
\node[below=4pt] at (0,-2) {$\Delta$};
\draw (3,-2.08)--(3,-1.92);
\node[below=4pt] at (3,-2) {$M-1-A$};
\draw (6.3,-2.08)--(6.3,-1.92);
\node[below=4pt] at (6.3,-2) {$M-1$};
\draw[->,blue!60!black,thick] (.5,-.3)--(5.8,-1.68);
\draw[->,blue!60!black,thick] (2.8,-.3)--(3.5,-1.68);
\node[anchor=south,font=\footnotesize] at (1.65,.45) {possible surviving indices};
\node[align=center,anchor=south] at (3.15,.95) {A unit coefficient with $a\leq A$ forces minimum $t$-valuation $\leq A$.};
\node[anchor=south] at (3.2,1.5) {$c_a(z-i)^a\ \longmapsto\ c_at^a(z-i)^{M-1-a}$};
\end{tikzpicture}
\caption{Taylor reversal in the expansion argument (schematic, not to scale). For a normalized input with integral Taylor coefficients and a unit among $a\leq A$, every term surviving in the leading reduction of its image has $a\leq A$. The reduction therefore vanishes to order at least $M-1-A$ at $i$. The proof applies this to every vector in a $j$-dimensional subspace, with $A=a_{ij}$.}
\label{fig:taylor-reversal}
\end{figure}
